\subsection{Interpretability, Fairness, and Safety}

Evaluation also asks whether a model can be understood, audited, and constrained in the setting where it is used.

\begin{definition} \label{def:sl-interpretability} \hlt{Interpretability}\\
Interpretability methods assign human-usable explanations to model behavior.  Examples include feature attribution, counterfactual explanations, surrogate models, and mechanistic analysis of internal representations.
\end{definition}

\begin{definition} \label{def:sl-fairness-criteria} \hlt{Fairness Criteria}\\
Group fairness criteria compare error rates, acceptance rates, calibration, or other performance quantities across protected groups.  Different fairness criteria can conflict, so the mathematical condition must be tied to a policy goal.
\end{definition}

\begin{definition} \label{def:sl-safety-constraint} \hlt{Safety Constraint}\\
A safety constraint restricts model behavior outside ordinary average-case loss, for example by enforcing abstention, monotonicity, bounded action sets, human review, or worst-case risk limits.
\end{definition}

\begin{remark} \label{rem:sl-evaluation-beyond-accuracy} \hlt{Beyond Accuracy}\\
Interpretability, fairness, and safety are not single metrics.  They are collections of mathematical requirements, measurement protocols, and governance choices around a deployed learning system.
\end{remark}
